\documentclass[11pt]{article}
\usepackage[letterpaper,margin=1in,top=1.2in,bottom=1.15in]{geometry}
\usepackage[utf8]{inputenc}\usepackage[T1]{fontenc}\usepackage{lmodern}
\usepackage{amsmath,amssymb,amsthm}\usepackage{mathtools}\usepackage{graphicx}
\usepackage{enumitem}\usepackage{fancyhdr}\usepackage{microtype}\usepackage{booktabs}
\usepackage{lastpage}\usepackage{listings}\usepackage{caption}
\usepackage[colorlinks=true,linkcolor=blue,citecolor=blue,urlcolor=blue]{hyperref}
\usepackage{url}
\lstset{basicstyle=\ttfamily\scriptsize,breaklines=true,frame=single,columns=flexible}
\captionsetup{font=small,labelfont=bf}

\theoremstyle{plain}
\newtheorem{theorem}{Theorem}[section]
\newtheorem{proposition}[theorem]{Proposition}
\newtheorem{corollary}[theorem]{Corollary}
\newtheorem{lemma}[theorem]{Lemma}
\theoremstyle{definition}
\newtheorem{definition}[theorem]{Definition}
\newtheorem{example}[theorem]{Example}
\theoremstyle{remark}
\newtheorem{remark}[theorem]{Remark}
\newtheorem{observation}[theorem]{Observation}

\newcommand{\tg}[1]{\langle #1\rangle}
\newcommand{\pow}{\mathcal{P}}
\newcommand{\Conop}{\operatorname{Con}}
\newcommand{\Np}{\mathbf{N}^{+}}

\pagestyle{fancy}\fancyhf{}
\lhead{\footnotesize Premise Selection over ZFC by Learning to Rank}
\rhead{\footnotesize Tenneson \quad$\vert$\quad Version 4.1 \quad$\vert$\quad July 28, 2026}
\lfoot{\footnotesize Copyright \textcopyright\ 2026 Brian Tenneson. All Rights Reserved.}
\cfoot{\footnotesize btenneson2301.substack.com}
\rfoot{\footnotesize Page \thepage\ of \pageref{LastPage}}
\renewcommand{\headrulewidth}{0.4pt}\renewcommand{\footrulewidth}{0.4pt}
\setlength{\headheight}{14pt}\setlength{\footskip}{28pt}

\title{\vspace{-1cm}\textbf{\Large Premise Selection over ZFC by Learning to Rank}\\[0.5em]
{\large Search Effort as the Governing Metric,\\ with Case Studies in Metamath's \texttt{set.mm}}}
\author{Brian Tenneson\\ M.S., Applied Data Science\\ Currently Unaffiliated\\
\href{mailto:btenneson2301@baypath.edu}{btenneson2301@baypath.edu}}
\date{Version 4.1 \quad$\bullet$\quad July 28, 2026}

\begin{document}
\maketitle
\thispagestyle{empty}

\begin{abstract}
\noindent
An automated theorem prover working inside a large formal library spends most of
its effort deciding which of tens of thousands of established results bear on the
goal in front of it. We treat that decision as a ranking problem over Metamath's
\texttt{set.mm}, a formalisation of ZFC containing on the order of $43{,}900$
proved theorems, and make three contributions.

First, a metric. Premise selectors are conventionally scored by recall at a
fixed cutoff, which measures the head of the ranking and is silent about its
tail. We argue that the governing quantity for a prover is \emph{search effort}:
the fraction of the candidate pool that must be examined before every premise of
a proof has been found, since that is the point at which search may stop. Recall
and effort are independent --- a ranker can excel at one and fail at the other
--- so both are reported throughout, against two baselines that a learned model
must beat to justify itself.

Second, a method. We fit a linear scorer by a pairwise ranking loss over
candidate pairs drawn from a single goal, rather than by classifying
(goal, candidate) pairs independently. Since every reported quantity depends
only on the order induced within one goal, this aligns the training objective
with the measurement. It also has a structural consequence, isolated as
Proposition~\ref{prop:cancel}: features that are constant across a goal's
candidates cancel identically and cannot receive weight, so a class of silently
uninformative feature is excluded by construction rather than by inspection.

Third, evidence. On the first $8{,}000$ statements of \texttt{set.mm} under a
temporal split, the ranker attains recall@$10$ of $0.232$ and recall@$50$ of
$0.491$ against $0.174$ and $0.233$ for the strongest baseline, and reduces
search effort from $0.926$ to $0.387$ --- from reading $2{,}779$ of $3{,}000$
candidates to reading $1{,}161$. Four worked case studies show where the gain
comes from and where it does not: premises sharing notation with the goal are
retrieved at ranks $1$ to $6$ out of several thousand, while premises stated in
notation absent from the goal are retrieved in the hundreds. We identify that
asymmetry as the principal obstacle to further progress and propose co-citation
over the premise graph as the natural remedy.
\end{abstract}

\medskip
\noindent\textbf{Keywords:} premise selection, learning to rank, automated
theorem proving, Metamath, ZFC, search effort, proof libraries, RankNet,
co-citation.

\newpage
\tableofcontents
\newpage

\section{Introduction}\label{sec:intro}

A theorem prover operating on a small axiom set can afford to search
exhaustively. One operating inside a formal library cannot. Modern libraries ---
Metamath's \texttt{set.mm}, Isabelle's Archive of Formal Proofs, Lean's Mathlib
--- contain tens or hundreds of thousands of established results, and the
dominant cost of proving a new statement is not applying an inference rule but
deciding \emph{which} of those results to apply it to. This is the premise
selection problem, and it is why interactive provers ship with tools such as
Sledgehammer \cite{sledgehammer} whose function is precisely to guess a short
list of relevant lemmas.

This paper studies premise selection over ZFC as formalised in \texttt{set.mm}
\cite{megill}. The choice of corpus is deliberate and is discussed in
Section~\ref{sec:corpus}: a Metamath proof is literally a list of label
references, so the premise relation is recorded in the file rather than inferred,
and no proof reconstruction or natural-language processing intervenes between the
data and the learning problem.

\subsection{Contributions}

\begin{enumerate}[leftmargin=1.5em,itemsep=3pt]
\item \textbf{Search effort as the governing metric}
(Section~\ref{sec:metrics}). We define effort as the depth to which a ranking
must be read before every true premise has been seen, argue that this rather
than recall is what a prover experiences, and show that the two quantities
diverge sharply in practice.

\item \textbf{A ranking formulation} (Section~\ref{sec:method}). We fit the
scorer by a pairwise loss over candidates of a common goal.
Proposition~\ref{prop:cancel} records a structural benefit: goal-constant
features cannot absorb weight.

\item \textbf{Typecode separation as a preprocessing requirement}
(Section~\ref{sec:corpus}). In Metamath, references to syntax constructors
outnumber references to mathematics by a factor approaching one in two. We
quantify this and show that failing to separate them changes both what the data
look like and what a model learns.

\item \textbf{Empirical results with worked case studies}
(Sections~\ref{sec:results}--\ref{sec:cases}). Aggregate figures on
\texttt{set.mm} are accompanied by four theorems examined premise by premise,
which localise the method's strength and its failure mode more precisely than
averages can.
\end{enumerate}

\section{The corpus and the premise relation}\label{sec:corpus}

\subsection{Metamath}

A Metamath database is a sequence of statements. Each is introduced by a label
and asserts a string of tokens, the first of which is its \emph{typecode}. Two
kinds concern us. A statement whose typecode is \verb!|-! asserts provability; one
whose typecode is \texttt{wff}, \texttt{class} or \texttt{setvar} asserts that a
string is well formed. Axioms are marked \texttt{\$a} and proved statements
\texttt{\$p}, and a proved statement carries its proof as a list of the labels it
invokes.

\begin{definition}[Premise relation]\label{def:premise}
Let $C=(t_0,\dots,t_{N-1})$ be a Metamath database in file order. For a proved
statement $t_i$ let $R(t_i)$ be the set of labels its proof references. The
\emph{premises} of $t_i$ are
\[
P(t_i)\;=\;\{\,t_j \;:\; j<i,\ \text{label}(t_j)\in R(t_i),\ \text{typecode}(t_j)=\ \text{\texttt{|-}}\,\}.
\]
The remaining references, those to statements of syntactic typecode, are the
\emph{syntactic references} of $t_i$.
\end{definition}

\subsection{Why the typecodes must be separated}

Definition~\ref{def:premise} restricts to logical typecode, and this restriction
is not cosmetic.

\begin{observation}[Syntactic references dominate]\label{obs:syntax}
In the first $8{,}000$ statements of \texttt{set.mm}, of $83{,}923$ total
references $35{,}902$ --- $42.8\%$ --- are to statements of syntactic typecode.
Counting them as premises raises the mean number of premises per proof from
$6.2$ to $10.8$.
\end{observation}

The reason is structural rather than incidental. Metamath requires a proof to
construct its formulas before asserting anything about them, so the constructors
for implication, conjunction, membership and equality are invoked by essentially
every proof in the library, independently of its subject. Table~\ref{tab:cited}
makes the effect visible: computed over all references, the most-cited labels are
uniformly grammatical; computed over logical references only, they are the
inference lemmas one would expect.

\begin{table}[htbp]\centering
\small
\begin{tabular}{@{}llrllr@{}}
\toprule
\multicolumn{3}{c}{\textbf{all references}} & \multicolumn{3}{c}{\textbf{logical references only}}\\
\cmidrule(r){1-3}\cmidrule(l){4-6}
label & role & count & label & role & count\\
\midrule
\texttt{wcel} & syntax: $A\in B$ is a wff & 2946 & \texttt{syl} & syllogism & 775\\
\texttt{wa}   & syntax: conjunction & 2790 & \texttt{bitri} & biconditional transitivity & 641\\
\texttt{wceq} & syntax: equality & 2545 & \texttt{ax-mp} & modus ponens & 482\\
\texttt{cv}   & syntax: setvar as class & 2365 & \texttt{a1i} & antecedent introduction & 348\\
\texttt{wi}   & syntax: implication & 1814 & \texttt{ex} & exportation & 342\\
\texttt{wn}   & syntax: negation & 1314 & \texttt{vex} & setvar existence & 322\\
\bottomrule
\end{tabular}
\caption{Most-referenced labels in the first $8{,}000$ statements of
\texttt{set.mm}. Without typecode separation the premise relation is dominated by
grammar; a model fitted to it learns to predict punctuation.}
\label{tab:cited}
\end{table}

\begin{remark}[Generality]
The same distinction arises in any system separating syntax from assertion. In
Isabelle and Lean it is largely handled by the elaborator and is invisible in the
proof term; in Metamath it is explicit in the file, which is convenient for the
present purpose but requires the filter of Definition~\ref{def:premise}.
\end{remark}

\subsection{Corpus statistics}

After separation, the first $8{,}000$ statements comprise $7{,}899$ logical
statements and $101$ syntactic ones, with $48{,}021$ logical references. Proofs
cite a mean of $6.2$ premises, median $4$, maximum $101$; the distribution of
citations is heavily skewed, with a small number of inference lemmas appearing
in hundreds of proofs and the majority of statements cited once or twice.

\section{Metrics}\label{sec:metrics}

Fix a goal $g$ with true premise set $G$, and a candidate pool $\Pi$ of
statements available to it. A selector induces a total order $\rho$ on $\Pi$.

\begin{definition}[Recall]\label{def:recall}
$\displaystyle \mathrm{recall}@k \;=\; \frac{|G\cap\{\rho_1,\dots,\rho_k\}|}{|G|}$.
\end{definition}

\begin{definition}[Search effort]\label{def:effort}
$\displaystyle \mathrm{effort} \;=\; \frac{\max\{\,r : \rho_r\in G\,\}}{|\Pi|}$,
the depth to which the ranking must be read before every premise has been seen.
\end{definition}

\begin{remark}[Why effort is the governing quantity]\label{rem:whyeffort}
Recall@$k$ describes the head of a ranking. It is the right measure for a system
that will be handed a fixed-size shortlist and must do the best it can with it.
It is the wrong measure for a prover that needs \emph{all} the premises of a
proof, because such a prover cannot stop at $k$: if one premise lies at rank
$2{,}000$, a ranking that placed the other four at ranks $1$ to $4$ has saved it
nothing. Effort measures exactly the stopping point.

The two are independent. A ranking with four premises at ranks $1$--$4$ and the
fifth at $2{,}000$ has excellent recall@$5$ and effort near $1$; a ranking with
all five at ranks $40$--$50$ has poor recall@$5$ and good effort. Neither implies
the other, and Section~\ref{sec:results} exhibits both patterns in the same
experiment.
\end{remark}

\begin{definition}[Baselines]\label{def:baselines}
Two orderings serve as controls. \emph{Frequency} sorts candidates by how often
each has been cited in the training prefix. \emph{Chronological} sorts by
recency, most recently proved first; since a library is developed topically, a
goal's premises are frequently among the statements proved shortly before it,
and reading the library backwards is a strong and entirely unlearned strategy.
We treat chronological order as the brute-force benchmark.
\end{definition}

\begin{remark}[Both metrics against both baselines]
A single number invites a misleading summary in either direction. Reporting
effort alone flatters a selector that returns a short list and misses premises;
reporting recall alone flatters one that finds the easy premises and buries a
hard one. We therefore report both, against both baselines, throughout.
\end{remark}

\section{Method}\label{sec:method}

\subsection{Ranking rather than classification}

The natural first formulation treats premise selection as binary
classification: fit $\Pr(\text{premise}\mid g,c)$ over (goal, candidate) pairs
and sort candidates by the fitted probability. We adopt a ranking formulation
instead, for a reason visible in Definitions~\ref{def:recall}
and~\ref{def:effort}: every reported quantity is a function of the order induced
\emph{within a single goal}. Absolute probabilities are never compared across
goals and never consulted; only their relative order within a goal is.

\begin{definition}[Pairwise ranking loss]\label{def:ranknet}
Let $g$ be a goal, $p\in P(g)$ a premise, $n\notin P(g)$ a non-premise available
to $g$, and $x_p,x_n\in\mathbf{R}^d$ their feature vectors. The scorer is linear,
$s(x)=w\cdot x$. We fit $w$ by logistic regression on difference vectors,
\[
\sigma\bigl(w\cdot(x_p-x_n)\bigr)\to 1,
\qquad
\sigma\bigl(w\cdot(x_n-x_p)\bigr)\to 0,
\]
over all pairs $(p,n)$ drawn from a common goal, where
$\sigma(z)=(1+e^{-z})^{-1}$. This is RankNet \cite{burges} with a linear
scorer; the antisymmetric second family of examples prevents the objective from
being satisfied by a uniform inflation of scores.
\end{definition}

\begin{proposition}[Goal-constant features cannot receive weight]\label{prop:cancel}
Suppose a feature coordinate $j$ satisfies $x_c^{(j)}=\kappa_g$ for every
candidate $c$ available to a goal $g$, with $\kappa_g$ depending on $g$ alone.
Then the $j$th coordinate of every difference vector formed from $g$ is zero, the
partial derivative of the loss with respect to $w_j$ vanishes on every such
example, and $w_j$ is unaffected by the data.
\end{proposition}

\begin{proof}
For any pair $(p,n)$ of candidates of $g$ we have
$x_p^{(j)}-x_n^{(j)}=\kappa_g-\kappa_g=0$. The loss depends on $w$ only through
the inner products $w\cdot(x_p-x_n)$, whose derivative with respect to $w_j$ is
the $j$th coordinate of the difference vector, namely zero. Summing over
examples, $\partial\mathcal{L}/\partial w_j=0$ identically.
\end{proof}

\begin{remark}[Consequence for feature design]\label{rem:cancel}
Proposition~\ref{prop:cancel} is a modest observation with a practical use. A
feature depending only on the goal --- the goal's length, its depth in the
library, the number of its premises --- is uninformative for ranking, since it
shifts every candidate's score by the same amount. Under a classification
objective such a feature can nonetheless acquire a large coefficient, because it
genuinely predicts the class marginal, and will appear informative in a table of
learned weights. Under Definition~\ref{def:ranknet} it is excluded by
construction. The point generalises: the ranking formulation restricts the model
to exactly the information that can affect the reported metrics.
\end{remark}

\subsection{Features}\label{sec:features}

All features are properties of a (goal, candidate) pair and are computable from
the database without proof reconstruction. Writing $S_g,S_c$ for the token sets
of goal and candidate:

\begin{enumerate}[leftmargin=1.6em,itemsep=2pt]
\item \emph{Symbol overlap}, $|S_g\cap S_c|/|S_g\cup S_c|$, and the two
asymmetric variants $|S_g\cap S_c|/|S_c|$ and $|S_g\cap S_c|/|S_g|$.
\item \emph{Literal occurrence}: whether the candidate's statement occurs as a
substring of the goal's.
\item \emph{Citation frequency}: $\log(1+u_c)$ where $u_c$ counts citations of
$c$ in the reference prefix.
\item \emph{Recency}: the positional gap between goal and candidate.
\item \emph{Size}: the candidate's length, and its ratio to the goal's.
\item \emph{Rare-symbol sharing}: whether goal and candidate share a token in the
lowest quartile of corpus frequency.
\item \emph{Local co-citation}: how often the candidate is cited by the $40$
statements immediately preceding the goal.
\item \emph{Co-citation with anchors}: see Section~\ref{sec:cocite}.
\end{enumerate}

\subsection{Negative sampling}\label{sec:negatives}

Training pairs require non-premises, and their selection is not neutral. Uniform
sampling from the available prefix supplies negatives that are almost all
trivially irrelevant --- on a corpus of tens of thousands, a random statement
bears no relation to a given goal --- so a model trained against them learns to
reject the unrelated rather than to discriminate among the plausible, which is
the operative task.

We therefore draw a majority of negatives from the most-cited statements of the
reference prefix: lemmas that are plausible for any goal because they are used
everywhere, and which the model must learn to rank \emph{below} the goal's actual
premises.

\begin{remark}[Negatives must be selected orthogonally to the features]\label{rem:orthogonal}
A design constraint deserves emphasis, as it is easy to violate. If negatives are
sampled by a criterion that some feature also measures, that feature becomes a
detector for the sampling procedure. Concretely: drawing hard negatives from the
premises of the $40$ statements preceding the goal, while also supplying a
feature counting citations within that same window, makes the feature predict
membership in the negative class almost perfectly --- and a correctly fitted
model will assign it a large negative weight, degrading performance on genuine
data where no such relationship holds. Selecting negatives by global citation
frequency, a quantity no locality feature computes, avoids the coupling.
\end{remark}

\subsection{Protocol}\label{sec:protocol}

\begin{definition}[Temporal split]\label{def:split}
Given a corpus in derivation order and a fraction $p$, the first $\lfloor pN\rfloor$
statements supply training goals and the remainder supply test goals. A random
split would allow the model to train on statements postdating its test goals,
which does not correspond to any deployment condition.
\end{definition}

\begin{definition}[Reference prefix]\label{def:refcut}
Several features --- citation frequency, symbol rarity, co-citation --- are
defined relative to a corpus. The prefix over which these are computed, the
\emph{reference prefix}, is specified separately from the training cut and held
fixed across any experiment that varies the amount of training data, so that a
feature denotes the same quantity during fitting and scoring.
\end{definition}

\section{Results}\label{sec:results}

We report on the first $8{,}000$ statements of \texttt{set.mm}, temporal split at
$p=0.8$, $200$ held-out goals, candidate pools of $3{,}000$.

\begin{table}[htbp]\centering
\small
\begin{tabular}{@{}lrrrr@{}}
\toprule
& \textbf{ranker} & frequency & chronological & fused \\
\midrule
recall@1  & 0.051 & 0.036 & \textbf{0.085} & 0.053 \\
recall@5  & \textbf{0.155} & 0.092 & 0.146 & 0.123 \\
recall@10 & \textbf{0.232} & 0.131 & 0.174 & 0.168 \\
recall@50 & \textbf{0.491} & 0.279 & 0.233 & 0.476 \\
\midrule
effort & 0.387 & --- & 0.926 & \textbf{0.321} \\
candidates read (of 3000) & 1161 & --- & 2779 & \textbf{964} \\
\bottomrule
\end{tabular}
\caption{Aggregate performance on \texttt{set.mm}. The fused column is
reciprocal-rank fusion of the learned ranking with the frequency ordering.}
\label{tab:agg}
\end{table}

\begin{observation}[The baseline wins at the head]\label{obs:head}
Chronological order attains the best recall@$1$, at $0.085$ against the ranker's
$0.051$. A library's most recently proved statement is a strong single guess,
because development proceeds topically and a new theorem commonly uses the lemma
established immediately before it. No learned model in our experiments displaced
this baseline at $k=1$.
\end{observation}

\begin{observation}[The learned ranker wins everywhere else]\label{obs:tail}
From $k=5$ upward the ordering reverses and the margin widens: at $k=50$ the
ranker recovers $0.491$ of premises against $0.233$ for chronological order, a
factor of $2.1$. On effort the difference is starker still: the ranker requires
$1{,}161$ of $3{,}000$ candidates against $2{,}779$, and rank fusion with the
frequency ordering reduces this to $964$.
\end{observation}

\begin{remark}[Reading the two together]
Observations~\ref{obs:head} and \ref{obs:tail} are the empirical content of
Remark~\ref{rem:whyeffort}. Chronological order is better at its first guess and
much worse at completing the set. A system permitted one attempt should read the
library backwards; a system that must assemble every premise of a proof should
use the learned ranking, and the difference between them --- roughly $1{,}600$
candidates of reading --- is the practical contribution.
\end{remark}

\section{Case studies}\label{sec:cases}

Aggregate figures conceal structure. We examine four theorems of \texttt{set.mm}
premise by premise; in each the candidate pool is every logical statement
preceding the goal, several thousand in size.

\begin{table}[htbp]\centering
\small
\begin{tabular}{@{}llrrr@{}}
\toprule
theorem & statement & ranker & frequency & chronological\\
\midrule
\texttt{prcom} & $\{A,B\}=\{B,A\}$ & \textbf{0.080} & 0.978 & 0.404\\
\texttt{uncom} & $(A\cup B)=(B\cup A)$ & \textbf{0.077} & 1.000 & 0.932\\
\texttt{unass} & $((A\cup B)\cup C)=(A\cup(B\cup C))$ & \textbf{0.285} & 0.997 & 0.933\\
\texttt{elpri} & $A\in\{B,C\}\rightarrow(A=B\vee A=C)$ & \textbf{0.135} & 1.000 & 0.942\\
\bottomrule
\end{tabular}
\caption{Search effort on four theorems. Lower is better; the ranker improves on
the better baseline by factors of $3.3$ to $12.1$.}
\label{tab:cases}
\end{table}

\begin{example}[Commutativity of unordered pairs]\label{ex:prcom}
The statement $\{A,B\}=\{B,A\}$ has a three-premise proof: \texttt{df-pr}, that
$\{A,B\}=\{A\}\cup\{B\}$; \texttt{uncom}, that union commutes; and
\texttt{3eqtr4i}, chaining three equalities. The argument is
$\{A,B\}=\{A\}\cup\{B\}=\{B\}\cup\{A\}=\{B,A\}$.

Out of $4{,}643$ candidates the ranker places \texttt{df-pr} at rank $2$,
\texttt{3eqtr4i} at $150$ and \texttt{uncom} at $372$; frequency places them at
$4541$, $1036$ and $4067$. The definitional premise is found almost immediately,
because the goal contains the literal tokens $\{\,A\,,\,B\,\}=$ and so does
\texttt{df-pr}.
\end{example}

\begin{observation}[Where the method succeeds]\label{obs:succeed}
Across the four theorems, premises sharing notation with the goal are retrieved
at ranks $1$ to $6$: \texttt{uneqri} at $1$ and \texttt{elun} at $3$ for
\texttt{uncom}; \texttt{uneqri} at $4$ and \texttt{elun} at $5$ for
\texttt{unass}; \texttt{elprg} at $6$ for \texttt{elpri}; \texttt{df-pr} at $2$
for \texttt{prcom}. Symbol overlap is close to decisive in this regime.
\end{observation}

\begin{observation}[Where it fails]\label{obs:fail}
Two classes of premise are retrieved poorly. The first is generic inference
lemmas --- \texttt{orcom} at rank $312$ against frequency's $62$, \texttt{orbi1i}
at $933$ against $215$, \texttt{3bitrri} at $1163$ against $307$. These appear in
proofs across all of mathematics and share no vocabulary with any particular
goal, so citation frequency predicts them and overlap cannot.

The second is more fundamental. In Example~\ref{ex:prcom} the premise
\texttt{uncom} is ranked $372$: the proof routes through union, but the goal
$\{A,B\}=\{B,A\}$ contains no union symbol. \emph{Any} feature computed from
shared vocabulary is structurally incapable of retrieving a premise whose
notation does not appear in the goal.
\end{observation}

\begin{remark}[The two failure modes are different]
The first is addressable by combination: frequency and the learned ranker fail
on disjoint sets of premises, and Table~\ref{tab:agg} shows that fusing them
improves aggregate effort from $0.387$ to $0.321$. The second is not, because no
amount of reweighting recovers information the features never encoded.
\end{remark}

\section{Co-citation}\label{sec:cocite}

Observation~\ref{obs:fail} identifies the obstacle: a premise stated in notation
absent from the goal is invisible to vocabulary-based features. We propose a
feature that measures a lemma's company rather than its vocabulary.

\begin{definition}[Co-citation]\label{def:cocite}
For statements $a,b$ let $c(a,b)$ be the number of proofs in the reference prefix
citing both. Given a goal, let $A$ be a small set of \emph{anchors} --- candidates
the vocabulary features already rank highly --- and define the co-citation score
of a candidate $u$ as $\sum_{a\in A}c(u,a)$.
\end{definition}

\begin{remark}[Why this can reach what overlap cannot]
Return to Example~\ref{ex:prcom}. The goal's notation points unambiguously at
\texttt{df-pr}, which the ranker places at position $2$; \texttt{df-pr} therefore
serves as an anchor. If \texttt{uncom} is cited alongside \texttt{df-pr}
throughout \texttt{set.mm} --- as it is, since unfolding a pair into a union is
the standard route to any statement about pairs --- then \texttt{uncom} receives a
high co-citation score despite sharing no symbol with the goal. The mechanism is
that of collaborative filtering: relatedness is inferred from co-occurrence
rather than from content.
\end{remark}

\section{Related work}\label{sec:related}

Premise selection has been studied principally in the Isabelle and Mizar
ecosystems, where Sledgehammer's relevance filter \cite{sledgehammer} established
the practical value of restricting a goal's premise set before invoking an
external prover. Learning-based selectors followed, and surveys of the modern
neural literature are available \cite{survey}. Our contribution is not a stronger
model --- the scorer here is linear, and deliberately so --- but a reframing of
the evaluation and of the training objective, together with an analysis
localising where surface features suffice and where they cannot.

The pairwise ranking formulation is RankNet \cite{burges}, standard in
information retrieval and, so far as we are aware, less commonly applied to
premise selection than pointwise classification. The metric we advocate, search
effort, plays the role that rank-of-last-relevant-item plays in retrieval
settings where a user must find every relevant document rather than any.

\section{Limitations and further work}\label{sec:limits}

\paragraph{Scale.} Results are reported on the first $8{,}000$ statements of
\texttt{set.mm}, roughly a fifth of the corpus. Whether the margin over the
baselines is preserved, widened or eroded as the library grows is open, and the
question is sharpened by the heterogeneity of \texttt{set.mm}: its early portion
is propositional logic and its later portion set theory, and a single linear
scorer may not serve both.

\paragraph{Model class.} The scorer is linear. A tree ensemble can represent
interactions between features --- that vocabulary overlap should count for more
on a content lemma than on a generic inference rule --- which a linear model
cannot. We regard this as the most promising direction for improving the head of
the ranking, where Observation~\ref{obs:head} shows a baseline still winning.

\paragraph{Structural features.} All features here are computed from token sets.
Genuine structural matching --- testing whether a candidate's conclusion unifies
with a subterm of the goal modulo substitution --- would address
Observation~\ref{obs:fail} directly rather than circumstantially, at
substantially greater computational cost.

\paragraph{Harder targets.} The theorems examined in Section~\ref{sec:cases} have
shallow proofs that unfold a definition and apply a propositional lemma. Cantor's
theorem, that no function maps a set onto its power set, is the natural next
test: its proof constructs the diagonal set $\{y\in x : y\notin f(y)\}$ and
derives a contradiction, so its premises concern separation, function application
and negation --- machinery whose notation need not appear in the statement at
all. It is precisely the regime of Observation~\ref{obs:fail}.

\section{Conclusion}

Premise selection is a ranking problem, and it is worth measuring and training as
one. Measuring it as one means reporting search effort alongside recall, since a
prover that must assemble every premise of a proof is governed by the depth of
its worst-ranked premise and not by the quality of its top few guesses. Training
it as one means fitting the scorer to the order induced within a goal, which
matches the objective to the measurement and, by Proposition~\ref{prop:cancel},
excludes a class of uninformative feature by construction.

On ZFC as formalised in \texttt{set.mm} the resulting selector reads $1{,}161$ of
$3{,}000$ candidates where reading the library in order requires $2{,}779$, and
retrieves premises sharing notation with the goal at ranks $1$ to $6$ out of
several thousand. It does not retrieve premises stated in unfamiliar notation,
and the case studies show why: no feature computed from shared vocabulary can.
Whether co-citation over the premise graph closes that gap is the question we
regard as most worth answering next.

\section*{Reproducibility}
\addcontentsline{toc}{section}{Reproducibility}
All results were produced by a single self-contained Python program depending on
no third-party libraries. Every run writes a timestamped directory containing the
fitted model, the measured quantities, and the corpus statistics used.
\begin{lstlisting}
python predator4.py fetch set.mm
python predator4.py stats  --db set.mm --limit 8000
python predator4.py train  --db set.mm --limit 8000 -p 0.8
python predator4.py prove  --db set.mm --label prcom --limit 12000
\end{lstlisting}

\section*{Acknowledgements}
\addcontentsline{toc}{section}{Acknowledgements}
The author thanks AI assistants used as engineering and verification tools during
the preparation of this work, including for code review, verification of
cross-references, and the identification of several methodological faults during
development. All mathematical claims, interpretations and final decisions remain
the responsibility of the author.

\begin{thebibliography}{9}
\bibitem{burges} Burges, C., Shaked, T., Renshaw, E., Lazier, A., Deeds, M.,
Hamilton, N., \& Hullender, G. (2005). Learning to rank using gradient descent.
\textit{Proceedings of the 22nd International Conference on Machine Learning},
89--96.
\bibitem{megill} Megill, N. D., \& Wheeler, D. A. (2019). \textit{Metamath: A
Computer Language for Mathematical Proofs} (2nd ed.). Lulu Press.
\bibitem{sledgehammer} Paulson, L. C., \& Blanchette, J. C. (2010). Three years
of experience with Sledgehammer, a practical link between automatic and
interactive theorem provers. \textit{Proceedings of IWIL}.
\bibitem{survey} Li, Z., Sun, J., Murphy, L., Su, Q., Li, Z., Zhang, X., Yang,
K., \& Si, X. (2024). A survey on deep learning for theorem proving.
\textit{Conference on Language Modeling}.
\bibitem{tenneson} Tenneson, B. (2026). \textit{Depths of a Simulation and
Formalized Self-Awareness}, Merged Edition.
\end{thebibliography}

\end{document}
